%! Logarithmic Functions — Unit 2, Lesson 2.11 — Warm-Up (Answer Key)
\documentclass[10pt]{article}
\usepackage{precalculus-article}
\usepackage{precalculus-key}   % pulls in -boxes; do NOT also load -boxes
\newcommand{\TallMath}[1]{$\displaystyle #1\rule[-1.4em]{0pt}{3.2em}$}

\begin{document}

\pageheader{Unit 2, Lesson 2.11}{Warm-Up --- Answer Key}
\namedateperiod

\vspace{0.12in}

\begin{enumerate}[leftmargin=1.8em, itemsep=14pt]

\item Evaluate each logarithm.

\vspace{4pt}
$\log_2 16 = $ \ans{4} \qquad $\log_3 1 = $ \ans{0} \qquad $\log_5 5 = $ \ans{1}

\begin{teachernote}
$\log_3 1 = 0$ because $3^0 = 1$. Students who write 1 are confusing the base with the argument.
$\log_5 5 = 1$ because $5^1 = 5$. Students who write 5 are not reading the log definition.
\end{teachernote}

\vspace{6pt}

\item Complete the table for $g(x) = 3^x$. Then answer the questions about its inverse
      $f(x) = \log_3 x$.

\vspace{6pt}
\begin{center}
\renewcommand{\arraystretch}{1.8}
\begin{tabular}{|c|c|}
\hline
\rowcolor{sky}
$x$ & $g(x) = 3^x$ \\
\hline
$-2$ & \ans{$1/9$} \\
\hline
$-1$ & \ans{$1/3$} \\
\hline
$0$  & \ans{$1$} \\
\hline
$1$  & \ans{$3$} \\
\hline
$2$  & \ans{$9$} \\
\hline
\end{tabular}
\end{center}

\vspace{6pt}
Which of the following are valid inputs for $f(x) = \log_3 x$?
Circle \textbf{all} that apply.

\vspace{6pt}
\begin{center}
$x = -3$ \qquad $x = 0$ \qquad \textcolor{keyred}{\textbf{$x = \tfrac{1}{9}$}} \qquad
\textcolor{keyred}{\textbf{$x = 1$}} \qquad \textcolor{keyred}{\textbf{$x = 27$}}
\end{center}

\begin{teachernote}
$x = -3$ and $x = 0$ are invalid: the domain of a logarithm is $x > 0$.
$x = 1/9$, $x = 1$, and $x = 27$ are all positive, so they are valid inputs.
\end{teachernote}

\vspace{6pt}

\item \textbf{True or False.} The function $f(x) = 2\log_5 x$ has a $y$-intercept.

\vspace{4pt}
Answer: \ans{False}

\vspace{4pt}
Brief justification: \ansline{The domain of $f$ is $x > 0$, so $x = 0$ is not in the domain.}
\ansline{There is no $y$-intercept; the graph has a vertical asymptote at $x = 0$.}

\end{enumerate}

\end{document}
